% =============================================================================
% Section 2: Background — Prior Work
% =============================================================================
% PURPOSE:
%   Summarise the contributions of Papers 1--8 in a unified narrative.
%   Paper 1 (brown note): airborne vs.\ mechanical coupling disparity.
%   Paper 7 (watermelon): acoustic ripeness inversion via shell resonance.
%   Paper 8 (Kac): identifiability lifting by asphericity.
%   Papers 2--6: supporting evidence (gas pockets, scaling laws, bladder,
%   borborygmi, sub-bass coupling).
% =============================================================================

\section{Background}
\label{sec:background}

\subsection{Shell-model overview}

The companion papers share a common forward model: a fluid-filled
viscoelastic shell whose undeformed midsurface is represented, to leading
order, as an oblate spheroid with semi-major axis $a$ and semi-minor axis
$c$.  The modelling strategy follows the classical shell-acoustics tradition
from Lamb's treatment of elastic-shell modal families~\cite{lamb1882},
through Love--Kirchhoff shell theory~\cite{Love1927}, to later
fluid--structure formulations in which structural inertia, fluid loading, and
curvature appear on equal footing~\cite{Junger1986}.  The abdominal formulation
in~\cite{BrowntoneP1} adapted that framework and used a Rayleigh--Ritz
approximation to retain the features that matter for the present programme:
curvature, wall compliance, internal loading, and mode shape.  A self-contained summary of the Ritz basis, boundary conditions, and
convergence evidence is provided in Appendix~\ref{sec:appendix_ritz}.
The capstone uses the resulting structure as
the common language in which the subsequent applications can be compared.

In schematic form, if $\mathbf{q}$ denotes the vector of Ritz amplitudes, the
kinetic and potential energies are written as
\begin{equation}
T = \frac{1}{2}\dot{\mathbf{q}}^{\mathsf T}
\left(\mathbf{M}_{s} + \mathbf{M}_{f}\right)\dot{\mathbf{q}},
\qquad
V = \frac{1}{2}\mathbf{q}^{\mathsf T}\mathbf{K}\mathbf{q},
\label{eq:background-energies}
\end{equation}
where $\mathbf{M}_{s}$ is the shell inertia matrix, $\mathbf{M}_{f}$ is the
added-mass contribution from the entrained fluid, and $\mathbf{K}$ is the
effective stiffness matrix arising from membrane strain, bending, and any
prestress retained in the model.  Stationarity of the Rayleigh quotient then
leads to the generalised eigenvalue problem
\begin{equation}
\left[\mathbf{K} - \omega^{2}\left(\mathbf{M}_{s} + \mathbf{M}_{f}\right)\right]
\mathbf{q} = 0,
\label{eq:background-eigenproblem}
\end{equation}
where $\omega$ is the angular natural frequency.  For nearly spherical shells
this structure reduces to the classical modal families identified by
Lamb~\cite{lamb1882}; for the more general fluid-loaded case, it fits squarely
within the coupled shell--acoustic framework described by Junger and
Feit~\cite{Junger1986}.  The important point for the present paper is that
geometry enters twice.  It shapes the admissible deformation fields through the
basis itself, and it enters the energetic balance through curvature-dependent
stiffness and fluid coupling.  The same geometric ingredients later control
both forcing efficiency and inverse conditioning.

\subsection{Companion studies: Stage I --- the question}

\emph{Stage~I} asked whether the brown-note claim contained any
physically defensible content, and the resulting analysis proved more informative
than the original claim itself.  The abdominal model in~\cite{BrowntoneP1},
formulated as a fluid-filled oblate shell, possesses a genuine low-frequency
flexural resonance.  The equivalent-sphere analytical estimate places the
second flexural mode at $f_{2}\approx\SI{3.95}{\hertz}$; the oblate Ritz model
at the canonical geometry ($a=\SI{0.18}{\metre}$, $c=\SI{0.12}{\metre}$) gives
the more accurate value
\begin{equation}
f_{2} = \SI{3.80}{\hertz}
\label{eq:background-f2}
\end{equation}
(single-term basis; converged value \SI{3.68}{\hertz}, see
Appendix~\ref{sec:appendix_convergence}).
That result mattered because it separated two questions that are often
conflated: whether a resonance exists, and whether any realistic forcing path
can excite it appreciably.  The modal calculation answered the first question
in the affirmative.  The coupling analysis answered the second in the negative
for airborne infrasound.  Because the acoustic field is both long-wavelength
and weakly matched to the relevant flexural deformation, the airborne pathway
does little work on the shell.  Mechanical forcing through whole-body vibration
projects onto the same mode family far more effectively, producing a coupling
ratio of approximately $R \approx 3.3 \times 10^{4}$ between the mechanical and
airborne responses~\cite{BrowntoneP1}.  The resonance was genuine, but the
direct acoustic pathway was ineffective.

The first important refinement, reported in~\cite{BrowntoneP2}, considered
tissue-constrained gas pockets as local compliant inclusions capable of much
stronger pressure-to-motion transduction.  This did not contradict the
whole-cavity result of~\cite{BrowntoneP1}.  Rather, it refined the
interpretation.  If one seeks a global
flexural response of the abdominal shell, airborne excitation is ineffective.
If one asks whether smaller substructures embedded within the abdomen might
respond locally, gas pockets are more plausible candidates.  The programme
therefore began with a division that remains useful throughout the capstone:
global shell resonance organises the large-scale spectrum, but local resonant
substructures may dominate transduction when coupling efficiency is the
relevant physiological question.

\subsection{Stage II --- evidence across scales, organs, and pathways}

\emph{Stage~II} tested whether the same shell-theoretic logic survived changes in
scale, organ, and excitation pathway.  The dimensionless reformulation
in~\cite{BrowntoneP3} recast the abdominal resonance problem and derived scaling
laws that made the low-order
flexural behaviour comparable across species~\cite{BrowntoneP3}.  Its point was
not merely zoological.  The governing frequency scales can be
expressed through a compact set of geometric and mechanical groups, which shows
that the infrasonic shell problem is structured by similarity,
rather than by the contingencies of one anatomical example.

The urinary-bladder analysis in~\cite{BrowntoneP4} transferred the framework and thereby tested it
on an organ whose geometry changes substantially with physiological
state~\cite{BrowntoneP4}.  Fill-dependent resonance in the bladder
cannot be interpreted as a simple size effect.  Increasing volume
changes characteristic length scales, but it also alters wall tension and
effective stiffness.  The resulting modal trajectory therefore reflects a
competition between geometry and stiffening, not a monotonic geometric trend.
This was the first clear sign, outside the original abdominal problem, that a
single scalar reduction of shape is unlikely to preserve the physically
relevant sensitivities.

The borborygmi study in~\cite{BrowntoneP5} shifted attention from infrasonic
flexure to an audible internal phenomenon.  Its constrained-bubble model
produced characteristic frequencies in the range
\SIrange{135}{440}{\hertz}, matching the clinical band of gut sounds.  The
mechanism was distinct from whole-cavity flexure, but the shell still mattered.
It set the confinement and loading conditions within which local gas pockets
oscillate.  The paper therefore broadened the programme without abandoning its
central theme: geometry continues to mediate what kind of spectrum is possible,
even when the relevant source is local rather than global.

The sub-bass analysis in~\cite{BrowntoneP6} returned to low frequencies and
extended the coupling question to the sub-bass regime, where concert exposures
and whole-body vibration are often discussed together despite acting through
rather different pathways~\cite{BrowntoneP6}.  The conclusion was consistent
with the abdominal analysis in~\cite{BrowntoneP1}.  Airborne
sub-bass remains a poor driver of substantial abdominal motion, whereas
structure-borne excitation can couple much more efficiently.  By the end of
Stage~II, the programme had assembled a coherent evidential base: across species,
organs, audible and infrasonic ranges, and airborne and mechanical forcing,
geometry continued to organise the modal spectrum while filtering how energy
enters it.

\subsection{Stage III --- from forward adequacy to inverse structure}

\emph{Stage~III} turned from forward resonance prediction to the inverse question
of what one can infer from a measured spectrum.  Watermelon tap-tone assessment
in~\cite{BrowntoneP7} provided the constructive example.  An equivalent-sphere
shell model can be forward-adequate for a
practical inversion task when the inferential target is narrow and the
geometric uncertainty is limited.  In that context the model need only map a
measured resonance to an effective rind stiffness with useful fidelity.  These
results therefore show that reduced geometry can be serviceable
when the application asks for classification or approximate property recovery,
rather than full geometric disentanglement.

The Kac-style identifiability analysis in~\cite{BrowntoneP8} exposed the
inverse cost of that simplification.
Asked in Kac-like language whether one can ``hear'' the shape of an organ, the
programme's answer became conditional rather than categorical.  If the shell is
reduced to a single equivalent radius, material and geometric sensitivities
become strongly correlated, and the inverse map is nearly singular.  The
condition-number analysis established that the equivalent-sphere model is
severely ill-conditioned, whereas the full oblate-shell model restores
practical separation of parameter effects.  Just as importantly,~\cite{BrowntoneP8}
revealed an oblate--prolate asymmetry: breaking spherical symmetry is not, by
itself, enough.  What matters is whether the perturbation changes how different
modes sample curvature.

A short theorem-oriented note on identifiability lifting~\cite{BrowntoneP9}
condensed that observation.  Its contribution was to show that
oblate deformations lift the near-spherical inverse degeneracy in a way that
prolate deformations do not improve comparably.  This matters to the capstone
because it rules out an imprecise summary.  One cannot simply say that
``asphericity helps''.  Some geometric departures from the sphere decorrelate
the spectral sensitivities that matter for inversion; others do not.  The
programme's theoretical turn therefore culminated in a sharper claim:
forward adequacy and inverse adequacy are distinct, and their separation is
governed by geometry.

\subsection{Broader context}

The Browntone series sits at the meeting point of several established
literatures.  One is classical shell acoustics, in which curvature, modal
families, and fluid loading have long been treated as inseparable elements of
the forward problem~\cite{Love1927,Leissa1973,Soedel2004,Junger1986}.  Another
is spectral inverse theory, from Kac's celebrated question and the
Gordon--Webb--Wolpert counterexample to engineering treatments of vibration
inversion and model updating~\cite{Kac1966,Gordon1992,Gladwell2004,Gladwell2005,
MottersheadFriswell1993,FriswellMottershead1995}.  A third is the practical
language of identifiability, familiar from inverse elasticity and biomedical
imaging, in which the issue is not merely uniqueness in principle but whether
the measured observables contain enough independent information to distinguish
the parameters of interest under realistic sensitivity and noise structures
~\cite{Raue2009,Doyley2012,parker2011imaging,Sack2023}.

What the present programme adds is a concrete family of case studies in which
these literatures can be read together.  The same geometric framework explains
why a soft organ can possess a low-frequency shell resonance yet remain almost
immune to direct airborne excitation; why local gas pockets and gut sounds are
better interpreted as shell-constrained substructures than as whole-cavity
breathing modes; why a reduced model can succeed as a forward surrogate for
fruit-ripeness assessment; and why that same reduction can fail badly when used
for multi-parameter inversion.  In biomedical and engineering settings alike,
the practical question is rarely whether resonance exists in the abstract.  It
is whether the observable spectrum contains geometry-specific information that
survives modelling reduction and experimental uncertainty.  The companion
papers suggest that this question is best answered not by separating shell
mechanics, acoustics, and inversion into different problems, but by treating
geometry as the common mediator of all three.
